\chapter{Real Structures on a Group}\label{chapter:realstructures}

The key object of our study are pairs $H\leq G$, where $G$ is a group and $H$ is an index 2 subgroup. In this chapter we consider such pairs from the point of view of $H$.

%\begin{defn}
A \emph{Real structure} on $H$ is an injective group homomorphism $\phi \colon H \rightarrow G$,
where $\phi(H) \leq {G}$ is an index two subgroup.
An \emph{isomorphism of Real structures} 
    \[
        \psi \colon (\phi_1 \colon H \rightarrow {G}_1) \rightarrow (\phi_2 \colon H \rightarrow {G}_2)
    \]
    is a group isomorphism $\psi \colon {G}_1 \rightarrow {G}_2$ which restricts to an inner automorphism of $H$.
    %\end{defn}
    We denote it by $(\phi_1 \colon H \rightarrow {G}_1) \simeq (\phi_2 \colon H \rightarrow {G}_2)$
    or simply ${G}_1 \simeq {G}_2$, when the maps are clear from the context. 
Our goal in this chapter is to classify Real structures on $H$ up to an isomorphism.    

We consider this definition of isomorphism, because inner automorphisms do not ``change'' representations of $H$, which is an important consideration for us. Instead of asking $\psi|_H$ to be an inner automorphism, however, a one might demand only that $\psi (H) = H$. We describe how the results from this chapter can also be used to understand this weaker notion of isomorphism in Section \ref{sect:weakequiv}.

\section{Elementary Approach}\label{sect:elementarydesc}
Let $H$ be a group. For $\alpha \in \Aut(H)$ and $h \in H$, consider a disjoint union of two copies of $H$
\[
G(\alpha, h) \coloneqq H \coprod Hw,
\]
where the second is written as $Hw$, using a formal symbol $w$. 
Define a binary operation $\cdot$ by
\[
\begin{array}{c|cc}
   \cdot & g_2 & g_2 w \\ \hline
   g_1   & g_1 g_2 & g_1 g_2 w \\
   g_1 w & g_1 \alpha(g_2) w & g_1 \alpha(g_2) h
\end{array}
\]
for all $g_1, g_2 \in H$. In particular, $w^2 = h$ and $w g = \alpha(g) w$ for any $g \in H$. Recall the adjoint map
\[
\Ad : H \rightarrow \Aut (H), \ \ \Ad (g): x \mapsto gxg^{-1} \, . 
\]
\begin{lemma}\label{lem:elementary1}
	The following statements hold for any group $H$. 
	\begin{itemize}
        \item $G(\alpha, h)$ defines a group if and only if
          \begin{equation} \label{Conditions_alpha_h}
            \Ad(h) = \alpha^2 \ \mbox{ and } \ \alpha(h) = h.
            \end{equation}
	\item If  $G(\alpha, h)$ is a group, $H \leq G(\alpha, h)$ is a Real structure on $H$.
	\item Suppose $H \leq {G}$ is a Real structure. Choose  $u \in {G} \setminus H$. Set
          $h \coloneqq u^2$ and $\alpha(h) \coloneqq \Ad(u)$. Then there is a canonical isomorphism
          \[ {G} \xrightarrow{\cong} G(\alpha, h), \ g\mapsto g, \ gu \mapsto gw\]
          of Real structures on $H$, which is the identity on $H$.
\end{itemize}
\end{lemma}

\begin{proof}
	Suppose first that $G(\alpha, h)$ defines a group. Then for any $g \in H$, by associativity
    \[
        (w \cdot w) \cdot g = w \cdot (w \cdot g)
    \]
    so $hg = \alpha^2(g)h$ and, hence, $\Ad(h) = \alpha^2$. Furthermore, $w$ has a left inverse $h^{-1} w$ and a right inverse $\alpha(h^{-1}) w$, so these must coincide and $\alpha(h) = h$.
    
    Conversely, suppose that $\Ad(h) = \alpha^2$ and $\alpha(h) = h$. The identity $e \in H$ provides an identity of $G(\alpha, h)$ because $\alpha(e) = e$. For associativity, suppose that $g_1 w^a, g_2 w^b, g_3 w^c \in G(\alpha, h)$. We want to show that
	\[
		g_1 w^a \cdot (g_2 w^b \cdot g_3 w^c) = (g_1 w^a \cdot g_2 w^b) \cdot g_3 w^c.
	\] 
	If $a = 0$ this is easy verified, so we may assume that $a = 1$. If $b = 0$, the the above identity becomes
	\[
	g_1 w \cdot (g_2 g_3 w^c) = (g_1 \alpha(g_2) w) \cdot g_3 w^c.
	\]
	which follows because $\alpha(g_2 g_3) = \alpha(g_2) \alpha(g_3)$. So we may also assume that $b = 1$. When $c= 0$ the identity becomes
	\[
	g_1 w \cdot (g_2 \alpha(g_3) w) = (g_1 \alpha(g_2) h) \cdot g_3
	\]
	or equivalently, simplifying both sides,
	\[
	g_1 \alpha(g_2) \alpha^2(g_3) h = g_1 \alpha(g_2) h g_3
	\]
	which follows from the fact that $\alpha^2(g_3) h = \Ad(h)(g_3) h = h g_3 h^{-1} h = h g_3$.
	
	When $c = 1$, the identity is
		\[
	g_1 w \cdot (g_2 \alpha(g_3) h) = (g_1 \alpha(g_2) h) \cdot g_3 w
	\]
	or equivalently, simplifying both sides,
	\[
	g_1 \alpha(g_2) \alpha^2(g_3) \alpha(h) w = g_1 \alpha(g_2) h g_3 w,
	\]
	which follows from the fact that $\alpha(h) = h$ and the same reasoning as above.
	
	For inverses, any element of $H$ has a two-sided inverse, so consider $g_1 w \in Hw$. An element $g_2 w \in Hw$ satisfies $(g_1 w)(g_2 w) = e$
        if and only if $g_1 \alpha(g_2) = h^{-1}$, so set $g_2 = \alpha^{-1}(g_1^{-1} h^{-1})$. Then $g_2 w$ defines a right inverse to $g_1 w$, which is also a left inverse because $G(\alpha, h)$ is a semigroup.
	
	The remaining statements are straightforward to verify.
\end{proof}
We can reinterpret this result in the following way. Pick any $\alpha \in \Aut(H)$ and $h \in H$ and
consider the following group, a quotient of the free product 
\[
G' (\alpha, h) \coloneqq \frac{H\ast \bZ}{N} = \frac{H \ast \langle w \rangle}{\langle \langle w^2=h, \  \ \forall g\in H  \ wg = \alpha (g) w \rangle \rangle} \; .
\]
This group comes with a natural map $H\rightarrow G' (\alpha, h), \ g\mapsto g$. 
\begin{cor}\label{cor:elementary} The following statements are equivalent.
\begin{enumerate}  
\item The natural map $H\rightarrow G' (\alpha, h)$ is a Real structure on $H$.
\item The natural map $H\rightarrow G' (\alpha, h)$ is injective.
\item The conditions~\eqref{Conditions_alpha_h} hold.
\end{enumerate} 
If these are satisfied, then the natural map $G'(\alpha, h) \rightarrow G(\alpha, h)$ is an isomorphism of Real structures.
\end{cor}

\begin{proof}
  Clearly (1) implies (2). To show that (2) implies (3), suppose that $H\rightarrow G' (\alpha, h)$ is injective. Then for any $g \in G$, $hgh^{-1} = w^2 g w^{-2} = \alpha^2(g)$ in $G'(\alpha, h)$, and thus by injectivity $\alpha^2 = \Ad(h)$. Similarly, $w$ commutes with $h$ in $G'(\alpha, h)$, and thus $hw = wh = \alpha(h)w$, so $h = \alpha(h)$.
  
  Finally, to see that (3) implies (1), note that if the conditions ~\eqref{Conditions_alpha_h} are satisfied then $G(\alpha, h)$ is a group by Lemma \ref{lem:elementary1}, and so from the presentation of $G'(\alpha,h)$ there is a surjective homomorphism $G'(\alpha,h) \rightarrow G(\alpha,h)$ which commutes with $\iota \colon H \rightarrow G'(\alpha,h)$ and $H \hookrightarrow G(\alpha,h)$. From the surjectivity, $[G'(\alpha,h) : \iota(H)] \geq [G(\alpha,h): H] = 2$, whereas from the presentation of $G'(\alpha,h)$, $[G'(\alpha,h) : \iota(H)] \leq 2$, so $[G'(\alpha,h) : \iota(H)] = 2$, which forces the surjection $G'(\alpha,h) \rightarrow G(\alpha,h)$ to be an isomorphism.
\end{proof}

Let us write $\fR(H)$ for the set of Real structures on $H$ up to isomorphism.
Lemma~\ref{lem:elementary1} yields the following description of this set. Set
\[
  \widetilde{\fR}(H) \coloneqq \{ (\alpha,h)\in \Aut(H)\times H \,\mid\, \mbox{equations }\eqref{Conditions_alpha_h}\mbox{ hold}\}.
\]
The Lemma~\ref{lem:elementary1} tells us that the map
\[
\widetilde{\fR}(H) \rightarrow \fR(H), \qquad (\alpha, h) \mapsto G(\alpha, h)
\]
is well-defined and surjective. It remains only to pinpoint when the Real structures $G(\alpha_1, h_1)$ and $G(\alpha_2, h_2)$ are isomorphic. 
We denote their corresponding elements $w$ by $w_1$ and $w_2$.
%\subsection{An Alternative Exposition}

\begin{lemma}\label{lem:modifiedelementary2}
  Suppose that $G(\alpha_1, h_1)$ and $G(\alpha_2, h_2)$ define groups, and $a, b \in H$. Then the formula
  \[
  \phi_{a,b}(g) = a^{-1} g a, \qquad \phi_{a,b}(g w_1) = (a^{-1} g a b) w_2,
  \]
  for $g \in G$ defines an isomorphism $\phi_{a,b} \colon G(\alpha_1, h_1) \rightarrow G(\alpha_2, h_2)$ of Real structures if and only if
  \begin{equation}\label{copyH2action}
  a b  \alpha_2(b) h_2 a^{-1} = h_1
  \ \mbox{ and } \ 
%  \Ad(a)^{-1} \Ad(b) \alpha_2 \Ad (a) = \alpha_1 \,
    \Ad(a b  \alpha_2(a)^{-1}) \alpha_2 = \alpha_1 \, ,
\end{equation}
and any isomorphism from $G(\alpha_1, h_1)$ to $G(\alpha_2, h_2)$ is of this form.
\end{lemma}
\begin{proof}
  Using Corollary \ref{cor:elementary}, such a formula defines a well-defined group homomorphism if and only if the defining relations of $G'(\alpha_1, h_1)$ are preserved by $\phi_{a,b}$. A routine calculation in $G(\alpha_2, h_2)$ yields that these are equivalent to the relations (\ref{copyH2action}) in $H$. It is clear that once $\phi_{a,b}$ is well-defined, it is an isomorphism.

  Finally, consider a isomorphism of Real structures $\phi: G(\alpha_1, h_1) \to G(\alpha_2, h_2)$.
  On $H$, $\phi$ is inner, which determines $a$, and $b$ is determined by $\phi (w_1) = bw_2$.
  Inevitably, $\phi = \phi_{a,b}$.
\end{proof}

Consider the semidirect $H \ltimes H$, where $H$ acts on itself via $\Ad \colon H \rightarrow \Aut(H)$.

\begin{cor}\label{cor:HHaction}
Suppose that $G(\alpha_1, h_1)$ and $G(\alpha_2, h_2)$ define groups. Then $G(\alpha_1, h_1) \cong G(\alpha_2, h_2)$ as Real structures if and only if there exist $a,b \in H$ which satisfy the relations (\ref{copyH2action}).

In particular, the group $H \ltimes H$ acts on the set $\widetilde{\fR}(H)$ on the left by the formula 
  \begin{equation}\label{copyHHaction_explicit}
     (a,b) : (\alpha , h ) \mapsto
    (\Ad(a b  \alpha(a)^{-1}) \alpha,  a b  \alpha(b) h a^{-1}) 
  \end{equation}
  so that the set of Real structures on $H$ can be described as a quotient
  $\fR(H) = (H \ltimes H) \backslash \widetilde{\fR}(H)$. 
\end{cor}

We can give a simpler description of $\fR(H)$ involving only one copy of $H$.

\begin{lemma}\label{lem:elementary3}
  Suppose that $G(\alpha_1, h_1)$ and $G(\alpha_2, h_2)$ define groups. Then $G(\alpha_1, h_1) \cong G(\alpha_2, h_2)$ as Real structures if and only if there is some $s \in H$ with
  \begin{equation}\label{copyH2action2}
  s  \alpha_2(s) h_2  = h_1
  \ \mbox{ and } \ 
    \Ad(s) \alpha_2 = \alpha_1 \, .
\end{equation} 
\end{lemma}
\begin{proof}
  Given such an $s$, we can define an isomorphism $G(\alpha_1, h_1) \rightarrow G(\alpha_2, h_2)$ using Lemma \ref{lem:modifiedelementary2} taking $a = 1$ and $b = s$. Conversely, given any isomorphism $\phi \colon G(\alpha_1, h_1) \xrightarrow{\sim} G(\alpha_2, h_2)$, if $\phi|_H = \Ad(t)$, then we obtain another isomorphism
  \[
  \psi \coloneqq \phi \circ \Ad(t^{-1}) \colon G(\alpha_1, h_1) \xrightarrow{\sim} G(\alpha_1, h_1) \xrightarrow{\sim} G(\alpha_2, h_2)
  \]
  which is the identity on $H$.
  Suppose $\psi(w_1) = s w_2$.  Then $\psi = \phi_{1,s}$ by Lemma \ref{lem:modifiedelementary2}, and therefore $s$ satisfies the required relations, again by Lemma \ref{lem:modifiedelementary2}.
\end{proof}

\begin{thm}\label{cor:Haction}
The group $H$ acts on the set $\widetilde{\fR}(H)$ on the left by the formula 
  \begin{equation}\label{Haction_explicit}
     s : (\alpha , h ) \mapsto
    (\Ad(s) \alpha,  s  \alpha(s) h) 
  \end{equation}
  so that the set of Real structures on $H$ can be described as a quotient
  $\fR(H) = H \backslash \widetilde{\fR}(H)$. 
\end{thm}

\begin{remark}
Corollary \ref{cor:HHaction} and Corollary \ref{cor:Haction} tell us that the $H \ltimes H$ orbits on $\tilde{\fR}(H)$ are the same as those of $H$. It is therefore natural to ask: given $(\alpha_1, h_1)$ and $(\alpha_2, h_2)$ in the same $H \ltimes H$-orbit, what is an explicit element of $H$ which maps $(\alpha_1, h_1)$ and $(\alpha_2, h_2)$?

Concretely, we have that
\begin{equation}\label{copycopyH2action}
  a b  \alpha_2(b) h_2 a^{-1} = h_1
  \ \mbox{ and } \ 
%  \Ad(a)^{-1} \Ad(b) \alpha_2 \Ad (a) = \alpha_1 \,
    \Ad(a b  \alpha_2(a)^{-1}) \alpha_2 = \alpha_1 \, ,
\end{equation}
hold, and we want to find $s \in H$ with
\begin{equation}\label{copycopyH2action2}
  s  \alpha_2(s) h_2  = h_1
  \ \mbox{ and } \ 
    \Ad(s) \alpha_2 = \alpha_1 \, .
\end{equation} 

It turns out that $s \coloneqq a b  \alpha_2(a)^{-1}$ does the trick. The second relation is clearly satisfied, and one can verify the first:
\begin{align*}
         s  \alpha_2(s) h_2  = a b  \alpha_2(a)^{-1}\alpha_2(a b  \alpha_2(a)^{-1}) h_2  & = 
         a b  \alpha_2(b)  \alpha^2_2(a^{-1}) h_2  = \\
         = a b  \alpha_2(b)  h_2 (a^{-1}) h_2^{-1} h_2 &
         = a b  \alpha_2(b)  h_2 a^{-1} = h_1.
  \end{align*}  
  \end{remark}

\begin{eg}
We can also see that the orbits of the trivial pair $(1,1)$ (corresponding to the trivial Real structure $G(1,1) = H \times C_2$) under both $H$ and $H \ltimes H$ are the same as follows. The map
\[
\mu \colon H \rightarrow \widetilde{\fR}(H), \quad x \mapsto ( \Ad (x),x^2)\, ,
\]
is $H \ltimes H$-equivariant, when $H$ is given the action of $H \ltimes H$ by $(a,b) \cdot x \coloneqq abxa^{-1}$. Indeed, 
\[
\mu((a,b) \cdot x) = \mu(abxa^{-1}) = (abxa^{-1}, a (bx)^2 a^{-1}),
\]
which is equal to
\[
(a, b) \cdot \mu(x) = (a,b) \cdot (\Ad (x),x^2) = (\Ad(ab\Ad(x)(a)^{-1})\Ad(x), ab \Ad(x)(b)x^2a^{-1}).
\]

The action of $H \ltimes H$ on $H$ is transitive, and therefore the image of $H$ under $\mu$ is exactly the $H \ltimes H$-orbit of $(1,1) = \mu(1)$. Therefore, we see that the $H$-orbit of $(1,1)$, which is, by definition, the image of $\mu$, agrees with the $H \ltimes H$ orbit of $(1,1)$.
\end{eg}

\section{Cohomological Approach}
Let $H$ be a group. In Section \ref{sect:elementarydesc} we classified Real structures on $H$ using an elementary approach. In this section we consider an alternative point of view and describe how Real structures on $H$ are classified via cohomology and the Eilenberg-MacLane Theorem \cite{EML}. As a reference we follow \cite[\S 1]{DJR}.
%Write $C_2 = \langle \epsilon \rangle$ for the group with two elements, where $\epsilon^2 = 1$.

We first note that a Real structure on $H$ is the same as an extension of $H$ by $C_2$.
Moreover, $(\phi_1 \colon H \rightarrow {G}_1) \simeq (\phi_2 \colon H \rightarrow {G}_2)$
%%are isomorphic
if and only if the corresponding extensions are \emph{inner virtually equivalent},
in the sense that there is an isomorphism $\psi \colon {G}_1 \rightarrow {G}_1$ such that the diagram
% https://q.uiver.app/#q=WzAsMTAsWzAsMCwiMSJdLFswLDEsIjEiXSxbMSwwLCJHIl0sWzEsMSwiRyJdLFsyLDAsIlxcd2lkZWhhdHtHXzF9Il0sWzIsMSwiXFx3aWRlaGF0e0dfMn0iXSxbMywxLCJDXzIiXSxbMywwLCJDXzIiXSxbNCwwLCIxIl0sWzQsMSwiMSJdLFsyLDMsIlxcaWQiLDJdLFswLDJdLFsxLDNdLFsyLDQsIlxccGhpXzEiXSxbMyw1LCJcXHBoaV8yIl0sWzYsOV0sWzcsOF0sWzcsNiwiXFxpZCIsMl0sWzUsNl0sWzQsN10sWzQsNSwiXFxwc2kiLDJdXQ==
\[\begin{tikzcd}
	1 & H & {{G}_1} & {C_2} & 1 \\
	1 & H & {{G}_2} & {C_2} & 1
	\arrow[from=1-1, to=1-2]
	\arrow["{\phi_1}", from=1-2, to=1-3]
	\arrow["\psi|_H"', from=1-2, to=2-2]
	\arrow[from=1-3, to=1-4]
	\arrow["\psi"', from=1-3, to=2-3]
	\arrow[from=1-4, to=1-5]
	\arrow["\id"', from=1-4, to=2-4]
	\arrow[from=2-1, to=2-2]
	\arrow["{\phi_2}", from=2-2, to=2-3]
	\arrow[from=2-3, to=2-4]
	\arrow[from=2-4, to=2-5]
\end{tikzcd}\]
commutes and $\psi|_H$ is an inner automorphism. We call these extensions {\em equivalent} if additionally $\psi|_H=\id$. Observe that, as in the proof of Lemma~\ref{lem:elementary3}, two extensions are equivalent if and only if they are inner virtually equivalent.

Let $\Out (H)$ denote the group of outer automorphisms of $H$, the quotient of $\Aut(H)$ by the subgroup of inner automorphisms.

\begin{defn}
  The \emph{coupling} $\theta$ of a Real structure $\phi \colon H \rightarrow {G}$ is the outer automorphism $\theta$ of $H$
  induced by conjugation by any $g\in G\setminus H$.
\end{defn}

As $\theta$ is viewed as an outer automorphism, this is independent of the choice of odd element $g \in G \setminus H$. Furthermore, isomorphic Real structures have the same coupling.
Note that, because $g^2\in H$ for any $g\in G$, we have $\theta^2 = 1$ in $\Out(G)$. In particular, the data of the coupling $\theta$ is equivalent to the data of a group homomorphism $\Phi_{\theta} \colon C_2 \rightarrow \Out(G)$.

\begin{defn}
We write $\fR_{\theta}(H)$ for the set of Real structures on $H$ with a fixed coupling $\theta$.
\end{defn}
By definition, there is a disjoint union decomposition
\begin{equation} \label{RealStr:RTheta}
\fR(H) = \bigsqcup_{\theta \in \Out(H), \ \theta^2 =1} \fR_{\theta}(H) \, . 
\end{equation}

\begin{defn}
  For $\theta \in \Out(H)$, let $Z_{\theta} \coloneqq Z(H)$, with $\bZ[C_2]$-module structure $\epsilon * z = \alpha(z)$ for any lift $\alpha \in \Aut(H)$ of $\theta$.
\end{defn}

Note that this is independent of the choice of $\alpha$, and for any such $\alpha$, $\alpha^2$ is inner, so $\alpha^2$ is trivial on $Z$ and the $\bZ[C_2]$-module structure is well-defined. 

From $\theta \in \Out(H)$, we therefore have the $C_2$-module $Z_{\theta}$ and may consider the cohomology groups $\HH^2(C_2, Z_{\theta})$ and $\HH^3(C_2, Z_{\theta})$. In order to describe $\fR_{\theta}$ in terms of these groups, we need a particular element $[\theta] \in \HH^3(C_2, Z_{\theta})$ associated to $\theta$, defined as follows.

We follow \cite[\S 7]{EML}. Choose a representative $\alpha \in \Aut(H)$ for $\theta$. Define the normalised function
\[ \tilde{\alpha} \colon C_2 \rightarrow \Aut(H), \ \ \tilde{\alpha}(1) = 1, \ \ \tilde{\alpha}(\epsilon) = \alpha \, .\]
As $\theta^2=1 \in \Out(H)$, $\alpha^2$ is inner, so we may choose some $t \in H$ such that $\alpha^2 = \Ad(t)$. This extends to another normalised function
\[ t: C_2 \times C_2 \rightarrow H, \ \ t(\epsilon, \epsilon) = t, \mbox{ and } t(x,y) = 1 \mbox{ if } 1\in \{x,y\}.\]
Then one obtains a normalised cocycle $f_\theta \in \HZ^3 (C_2 , Z_{\theta})$ defined by
\[
    f_\theta(x,y,z) \coloneqq \tilde{\alpha}(x)(t(y,z))t(x,yz)t(xy,z)^{-1}t(x,y)^{-1}.
\]
%We set $[\theta]$ to be the induced element of $\HH^3(C_2, Z_{\theta})$.
\begin{defn}
We set $[\theta] \coloneqq [f_\theta]\in \HH^3(C_2, Z_{\theta})$. 
\end{defn}
From Proposition \ref{prop:C2cohomology} there is a canonical identification between $\HH^3(C_2, Z_{\theta})$ defined in terms of normalised cocycles and the group $Z_{\theta}^1 / (\epsilon - 1) Z_{\theta}$, defined by sending a normalised cocycle $f$ to $f(\epsilon, \epsilon, \epsilon)$. In our case, the associated element is 
\begin{equation}
z_\theta \coloneqq  f_\theta(\epsilon , \epsilon , \epsilon ) =  \alpha (t)t^{-1} \in Z_{\theta}^1 / (\epsilon - 1) Z_{\theta}.
\end{equation}

With this element $[\theta]$, we can now state the Eilenberg-MacLane Theorem.
\begin{thm}[{\cite[Prop.\ 1.4]{DJR}}]\label{thm:EMLv1}
    Let $\theta \in \Out(H)$ such that $\theta^2=1$.
    \begin{itemize}
        \item If $[\theta] \neq 0$ in $\HH^3(C_2, Z_{\theta})$, then $\fR_{\theta}(H) = \emptyset$.
        \item If $[\theta] = 0$ in $\HH^3(C_2, Z_{\theta})$, then $\fR_{\theta}(H)$ has a free transitive group action of $\HH^2(C_2, Z_{\theta})$.
    \end{itemize}
\end{thm}

\section{Unification}\label{sect:unification} % of both approaches}
In this section we reconcile the elementary and cohomological approaches to describing $\fR(H)$.

\subsection{Existence}

We first consider the question of whether $\fR_{\theta}(H) \neq \emptyset$ from each point of view.

From Lemma~\ref{lem:elementary1} we have that $\fR_{\theta}(H) \neq \emptyset$ if and only if there exist $\alpha \in \Aut(H)$ and $h \in H$ with $\overline{\alpha} = \theta$, $\Ad(h) = \alpha^2$ and $\alpha(h) = h$. 

We can show directly that this is consistent with Theorem~\ref{thm:EMLv1}. %\ref{thm:EMLv2}.

\begin{lemma}\label{lem:deg3cohom}
    Suppose that $\theta \in \Out(H)$. Then $[f_{\theta}] = 0$ in $\HH^3(Z_{\theta})$ if and only if there are $\alpha \in \Aut(H)$ and $h \in H$ with $\overline{\alpha} = \theta$, $\Ad(h) = \alpha^2$ and $\alpha(h) = h$.
\end{lemma}

\begin{proof}
  By construction, $[f_{\theta}] = 0$ in $\HH^3(C_2, Z_{\theta})$ if and only if $z_{\theta} = 1$ in $Z_{\theta}^1 / (\epsilon - 1) Z_{\theta}$.

  If so, let $\alpha \in \Aut(H)$ and $t \in H$ be the elements involved in the construction of $z_{\theta}$, so $\overline{\alpha} = \theta$, $\Ad(t) = \alpha^2$, and $z_{\theta} = \alpha(t)t^{-1}$. As $z_{\theta} = 1$, there is some $z \in Z_{\theta}$ with $\alpha(t)t^{-1} = \alpha(z) z^{-1}$, and so we may set $h \coloneqq tz^{-1}$ and compute
    \[
        \alpha(h)h^{-1} = \alpha(t)\alpha(z)^{-1}t^{-1} z = 1 \ 
        \mbox{ and } \ \Ad(h) = \Ad(t) = \alpha^2.
    \]

    Conversely, given $\alpha \in \Aut(H)$ and $h \in H$ with $\overline{\alpha} = \theta$, $\Ad(h) = \alpha^2$ and $\alpha(h) = h$, we may use these elements to construct the element $z_{\theta}$, the value of $z_{\theta} \in Z_{\theta}^1 / (\epsilon - 1) Z_{\theta}$ being independent of this choice. For this pair, $z_{\theta} = \alpha(h)h^{-1} = 1$ on the nose, and so in particular $z_{\theta} = 1$ in $Z_{\theta}^1 / (\epsilon - 1) Z_{\theta}$.
\end{proof}

\subsection{Group Action}\label{sect:freetransgroupaction}

The second part of Theorem \ref{thm:EMLv1} says that when $\fR_{\theta}(H)$ is non-empty, there is a free and transitive group action of $\HH^2(C_2, Z_{\theta})$ on $\fR_{\theta}(H)$. Let us now also show this from the elementary point of view, and the description of $\HH^2(C_2, Z_{\theta})$ as $Z^{C_2}_{\theta} / (\epsilon + 1)Z_{\theta}$ of Proposition \ref{prop:C2cohomology}.

Suppose that $\theta \in \Out(H)$ with $\fR_{\theta}(H) \neq \emptyset$, and pick an element of $\fR_{\theta}(H)$, which by Lemma \ref{lem:elementary1} we may choose to represented by some $G(\alpha, h)$ with $\alpha \in \Aut (H)$ and $h \in H$ such that $\alpha^2 = \Ad(h)$ and $\alpha(h) = h$.

We can then define an action of $Z_{\theta}^{C_2} = \{z \in Z_{\theta} \mid \alpha(z) = z\}$ on $\fR_{\theta}(H)$ as follows.

\begin{defn}
  For $z \in Z_{\theta}^{C_2}$, let $z*G(\alpha, h)\coloneqq G(\alpha, zh)$.
\end{defn}

This is well-defined because $\Ad(zh) = \Ad(h)$ and $\alpha(zh) = \alpha(z) \alpha(h) = zh$. This is furthermore independent of the choice of representative: if $G(\alpha_1, h_1) \simeq G(\alpha_2, h_2)$, then by Lemma~\ref{lem:elementary3} there is some $s \in H$ such that $\Ad(s) \circ \alpha_2 = \alpha_1$ and $h_1 = s \alpha_2(s) h_2$, and so $h_1 z = s \alpha_2(s) h_2 z$. Hence, $G(\alpha_1, h_1 z) \simeq G(\alpha_2, h_2 z)$ by Lemma~\ref{lem:elementary3}.

\begin{lemma}\label{lem:freetransaction}
    The construction above defines a transitive group action of $Z_{\theta}^{C_2}$ on $\fR_{\theta}(H)$ with all stabilisers equal to $\{\alpha(z)z \mid z \in Z_{\theta}\}$. In particular, the induced action of $\HH^2(C_2, Z_{\theta})$ on $\fR_{\theta}(H)$ is transitive and free.
\end{lemma}

\begin{proof}
  It is clear that this defines a group action. For transitivity, consider $G(\alpha_1, h_1),G(\alpha_2, h_2)\in \fR_{\theta}(G)$. As $G(\alpha_1, h_1)$ and $G(\alpha_2, h_2)$ have the same coupling, $\Ad(t) \circ \alpha_2 = \alpha_1$ for some $t \in H$, and therefore by Lemma~\ref{lem:elementary3},
    \[
    G(\alpha_2, h_2) \;\simeq\; G(\Ad(t) \circ \alpha_2, t \alpha_2(t)h_2) \;\simeq\; G(\alpha_1, t \alpha_2(t) h_2) \, .
    \]
    In particular, we may assume that $\alpha_1$ and $\alpha_2$ are both equal to some common $\alpha \in \Aut(H)$. 

    We have that $\Ad(h_1h_2^{-1}) = \Ad(h_1) \Ad(h_2)^{-1} = 1$, hence $h_2 = h_1 z$ for some $z \in Z$. Furthermore, $h_2 = \alpha(h_2) = \alpha(h_1) \alpha(z) = h_2 \alpha(z)$, so $\alpha(z) = z$, hence $z \in Z_{\theta}^{C_2}$ and $z * G(\alpha, h_1) = G(\alpha, h_2)$ in $\fR_{\theta}(H)$.

    To compute the stabilisers, suppose that $z \in Z_{\theta}^{C_2}$ and $G(\alpha, h) \in \fR_{\theta}(H)$. Then $z *  G(\alpha, h) \simeq G(\alpha, h)$ if and only if $G(\alpha, zh) \simeq G(\alpha, h)$, which holds if and only if there is some $s \in Z$ with $zh = s \alpha(s) h$, or equivalently, if and only if $z \in \{\alpha(z)z \mid z \in Z_{\theta}\}$.
\end{proof}

\section{Examples}

In this section we give examples of some families of groups for which we can use the theory developed above to describe all Real structures.

\subsection{Abelian Groups}\label{sect:realstrabgroups}

In this section we assume that $H$ is abelian. Then $\Out(H) = \Aut(H)$, and for any $\alpha \in \Aut(H)[2]$, we always have that $\fR_{\alpha}(H) \neq \emptyset$, as one can always explicitly take $h = 1$ in Lemma \ref{lem:deg3cohom}. Furthermore, $Z(H) = H$, and so there is a canonical identification
\[
  \frac{\{h \in H \mid \alpha(h) = h\}}{\{\alpha(z)z \mid z \in H\}} \xrightarrow{\sim} \fR_{\alpha}(H), \qquad h \mapsto G(\alpha, h).
\]

\begin{eg}
Suppose that $H = \bZ$ is infinite cyclic. Then $\Aut(\bZ)[2] = \{\pm 1\}$. If $\alpha$ is the identity, then 
\[
\frac{\{h \in \bZ \mid \alpha(h) = h\}}{\{\alpha(z) + z \mid z \in \bZ\}} = \bZ / 2 \bZ.
\]
The trivial element corresponds to $\bZ \leq \bZ \times \bZ / 2 \bZ$, and the non-trivial element to $\bZ \leq \frac{1}{2}\bZ$.

If instead $\alpha(1) = -1$, then $\{z \in \bZ \mid -z = z\} = 0$, so $|\fR_{\alpha}(H)| = 1$ and there is a unique Real structure on $\bZ$ which induces this automorphism on $\bZ$. It is the infinite dihedral group $\bZ \leq D_{\infty}$.
\end{eg}

Now let us consider the case of finite cyclic groups $C_n$, for $n \geq 1$. In what follows we use the canonical identification
\[
(\bZ / n \bZ)^{\times} \xrightarrow{\sim} \Aut(C_{n}), \qquad t \mapsto [\alpha_t \colon x \mapsto x^t].
\]
We first treat the case of cyclic groups of prime power order.

\begin{eg}\label{eg:RealStructuresOddPrimePower}
Suppose that $H = C_{p^n} = \langle x \rangle$, where $n \geq 1$ and $p$ is an odd prime. From the isomorphism
\[
(\bZ / p^n \bZ)^{\times} \cong \bZ / (p-1) \bZ \times \bZ / p^{n-1} \bZ
\]
and the fact that $p$ is odd, we have that $|\Aut(C_{p^n})[2]| = 2$, and therefore $\Aut(C_{p^n})[2] = \{\alpha_1, \alpha_{-1}\}$. 

If $\alpha = \alpha_1$, then 
\[
\frac{\{h \in C_{p^n} \mid \alpha_1(h) = h\}}{\{\alpha_1(z)z \mid z \in C_{p^n}\}} = C_{p^n} / C_{p^n}^2 = 0,
\]
and therefore $|\fR_{\alpha_1}(C_{p^n})| = 1$, with unique Real structure $C_{p^n} \leq C_{p^n} \times C_2$ (which is isomorphic to the Real structure $C_{p^n} \leq C_{2 p^n}$). If instead $\alpha = \alpha_{-1}$, then
\[
\{h \in C_{p^n} \mid \alpha_{-1}(h) = z\} = \{h \in C_{p^n} \mid h^2 = 1\} = 0,
\]
and therefore $|\fR_{\alpha_{-1}}(C_{p^n})| = 1$, consisting of $C_{p^n} \leq D_{2p^n}$.
\end{eg}

There are more Real structures, however, when the prime $p$ is $2$.

\begin{eg}\label{eg:RealStructuresEvenPrimePower}
  Suppose that $H = C_{2^n} = \langle x \rangle$, where $n \geq 1$. In general, when $\alpha = \alpha_1$,
  \[
  \frac{\{h \in C_{2^n} \mid \alpha_1(h) = h\}}{\{\alpha_1(z)z \mid z \in C_{2^n}\}} = C_{2^n} / C_{2^n}^2,
  \]
  and therefore there are two Real structures, namely $G(\alpha_1, 1)$ and $G(\alpha_1, x)$, which are given by $C_{2^n} \leq C_{2^n} \times C_2$ and $C_{2^n} \leq C_{2^{n+1}}$ respectively.
  
  To determine $\Aut(C_{2^n})[2]$, there are now three cases, depending on $n$.
  
  $\bullet$ When $n = 1$, $\Aut(C_2) = 1$, and $\fR_{\alpha_1}(C_2)$ consists of the two Real structures $C_{2} \leq C_{2}\times C_2$ and $C_2 \leq C_4$ by the above.

  $\bullet$ When $n = 2$, $\Aut(C_4) = \{\pm 1\}$, and $\fR_{\alpha_1}(C_2)$ consists of the two Real structures $C_{4} \leq C_{4}\times C_2$ and $C_4 \leq C_8$ by the above. For $\fR_{\alpha_{-1}}(C_2)$, we have that
  \[
  \frac{\{h \in C_4 \mid \alpha_{-1}(h) = h\}}{\{\alpha_1(z)z \mid z \in C_4\}} = C_4^2 / 1,
  \]
  and so $|\fR_{\alpha_{-1}}(C_2)| = 2$, given by $C_4 \leq D_8$ and $C_4 \leq Q_8$.

  $\bullet$ When $n \geq 3$,
  \[
  (\bZ / 2^n \bZ)^{\times} \cong \bZ / 2 \bZ \times \bZ / 2^{n-2} \bZ,
  \]
  and so $|\Aut(C_{2^n})[2]| = 4$. The elements $\{\pm 1, 2^{n-1} \pm 1\}$ all square to $1$ in $(\bZ / 2^n \bZ)^{\times}$, and therefore $\Aut(C_{2^n})[2] = \{\alpha_t \mid t \in \{\pm 1, 2^{n-1} \pm 1\}\}$. 

  When $t = \pm 1$, 
  \begin{align*}
    \sR_{\alpha_1}(C_{2^n}) &= \{C_{2^n} \leq C_{2^n} \times C_2, \: C_{2^n} \leq C_{2^{n+1}}\}, \\
    \sR_{\alpha_1}(C_{2^n}) &= \{C_{2^n} \leq D_{2^{n+1}}, \: C_{2^n} \leq Q_{2^{n+1}}\},
  \end{align*}
  similarly to above, where $Q_{2^{n+1}}$ is the generalised quaternion group of order $2^{n+1}$. 
  
  When $t = 2^{n-1} + 1$, we have that 
  \[
  \frac{\{h \in C_{2^n} \mid \alpha_t(h) = h\}}{\{\alpha_t(z)z \mid z \in C_{2^n}\}} = C_{2^n}^2 / C_{2^n}^2 = 0
  \]
  and therefore $|\fR_{\alpha_t}(C_{2^n})| = 1$, given by 
\[
C_{2^n} = \langle x \rangle \leq \langle x , w \mid x^{2^n}, w^2, wxw^{-1} = x^{2^{n-1} + 1} \rangle \eqqcolon M_{2^{n+1}},
\]
the \emph{modular maximal-cyclic group} or \emph{modular group} of order $2^{n+1}$.

When $t = 2^{n-1} - 1$, we have that
\[
\frac{\{h \in C_{2^n} \mid \alpha_t(h) = h\}}{\{\alpha_t(z)z \mid z \in C_{2^n}\}} = C_{2^n}^{2^{n-1}} / C_{2^n}^{2^{n-1}} = 0
\]
and therefore $|\fR_{\alpha_t}(C_{2^n})| = 1$, given by 
\[
C_{2^n} = \langle x \rangle \leq \langle x , w \mid x^{2^n}, w^2, wxw^{-1} = x^{2^{n-1} - 1} \rangle \eqqcolon \SD_{2^{n+1}},
\]
the \emph{semidihedral group} of order $2^{n+1}$.
\end{eg}

We can use the prime power order cases above to describe the general case. We write
\[
\phi \colon \bZ / n\bZ \xrightarrow{\sim} \bZ / 2^r \bZ \times \bZ / p_1^{n_1} \bZ \times \cdots \times \bZ / p_k^{n_k} \bZ,
\]
for the canonical isomorphism induced by projection, where $n = 2^r p_1^{n_1} \cdots p_k^{n_k}$.

\begin{cor}
Suppose that $n = 2^r p_1^{n_1} \cdots p_k^{n_k}$ a prime factorisation, with $r \geq 0$, $k \geq 0$, and $n_i \geq 1$ for all $i$, and $C_n = \langle x \rangle$. Then for $t \in (\bZ / n\bZ)^{\times}[2]$, $\sR_{\alpha_t}(C_n)$ is as in one of the following cases.
\begin{enumerate}
  \item[\emph{($r = 0$)}] $|\Aut(C_n)[2]| = 2^k$, $|\fR_{\alpha_t}(C_n)| = 1$, and $\fR_{\alpha_t}(C_n)$ is described by
  \[
    \langle x , w \mid x^n, w^2, wxw^{-1} = x^t \rangle.
  \]
  \item[\emph{($r = 1$)}] $|\Aut(C_n)[2]| = 2^k$, $|\fR_{\alpha_t}(C_n)| = 2$, and $\fR_{\alpha_t}(C_n)$ is described by
  \[
    \langle x , w \mid x^n, w^2, wxw^{-1} = x^t \rangle, \qquad \langle x , w \mid x^n, w^2 = x^{p_1^{n_1} \cdots p_k^{n_k}}, wxw^{-1} = x^t \rangle.
  \]
  \item[\emph{($r = 2$)}] $|\Aut(C_n)[2]| = 2^{k+1}$, $|\fR_{\alpha_t}(C_n)| = 2$, and there are two cases:
  \begin{itemize}
    \item If $t \equiv 1 \!\! \mod 4$, $\fR_{\alpha_t}(C_n)$ is described by
    \[
    \langle x , w \mid x^n, w^2, wxw^{-1} = x^t \rangle, \qquad \langle x , w \mid x^n, w^2 = x^u, wxw^{-1} = x^t \rangle,
    \]
   where $\phi(u) = (1,0,0, ... , 0)$.
  \item If $t \equiv -1 \!\! \mod 4$, $\fR_{\alpha_t}(C_n)$ is described by
  \[
    \langle x , w \mid x^n, w^2, wxw^{-1} = x^t \rangle, \qquad \langle x , w \mid x^n, w^2 = x^u, wxw^{-1} = x^t \rangle,
  \]
  where $\phi(u) =  (2,0,0, ... , 0)$.
  \end{itemize}
  \item[\emph{($r \geq 3$)}] $|\Aut(C_n)[2]| = 2^{k+2}$, and there are three cases:
  \begin{itemize}
  \item If $t \equiv 1 \!\! \mod 2^r$, $|\fR_{\alpha_t}(C_n)| = 2$, and $\fR_{\alpha_t}(C_n)$ is described by
    \[
    \langle x , w \mid x^n, w^2, wxw^{-1} = x^t \rangle, \qquad \langle x , w \mid x^n, w^2 = x, wxw^{-1} = x^t \rangle,
  \]
  where $\phi(u) = (1,0,0, ... , 0)$.
  \item If $t \equiv -1 \!\! \mod 2^r$, $|\fR_{\alpha_t}(C_n)| = 2$, and $\fR_{\alpha_t}(C_n)$ is described by
    \[
    \langle x , w \mid x^n, w^2, wxw^{-1} = x^t \rangle, \qquad \langle x , w \mid x^n, w^2 = x^{2^{r-1}}, wxw^{-1} = x^t \rangle,
  \]
  where $\phi(u) = (2^{r-1},0,0, ... , 0)$.
  \item If $t \equiv 2^{r-1} \pm 1 \!\! \mod 2^r$, $|\fR_{\alpha_t}(C_n)| = 1$, and $\fR_{\alpha_t}(C_n)$ is described by
    \[
    \langle x , w \mid x^n, w^2, wxw^{-1} = x^t \rangle.
  \]
  \end{itemize}
\end{enumerate}
In particular, 
\[
|\fR(C_n)| = \twopartdef{2^{k+r}}{r \leq 2}{2^{k+2} + 2^{k+1}}{r \geq 3}.
\]
\end{cor}

\begin{proof}
  From the isomorphism
  \[
  C_n \xrightarrow{\sim} C_{2^r} \times C_{p_1^{n_1}} \cdots C_{p_k^{n_k}} 
  \]
  there is a square of isomorphisms
      % https://q.uiver.app/#q=WzAsNCxbMSwwLCIoXFxiWiAvIDJeciBcXGJaKV57XFx0aW1lc30gXFx0aW1lcyAoXFxiWiAvIHBfMV57bl8xfSBcXGJaKV57XFx0aW1lc30gXFx0aW1lcyBcXGNkb3RzIFxcdGltZXMgKFxcYlogLyBwX2tee25fa30gXFxiWilee1xcdGltZXN9Il0sWzAsMCwiKFxcYlogLyBuXFxiWilee1xcdGltZXN9Il0sWzAsMSwiXFxBdXQoQ19uKSJdLFsxLDEsIlxcQXV0KENfezJebn0pIFxcdGltZXMgXFxBdXQoQ197cF8xXntuXzF9fSkgXFx0aW1lcyBcXGNkb3RzIFxcdGltZXMgXFxBdXQoQ197cF9rXntuX2t9fSkgIl0sWzEsMCwiXFxzaW0iXSxbMSwyLCJcXHNpbSIsMl0sWzMsMiwiXFxzaW0iLDJdLFswLDMsIlxcc2ltIiwyXV0=
\[\begin{tikzcd}
	{(\bZ / n\bZ)^{\times}} & {(\bZ / 2^r \bZ)^{\times} \times (\bZ / p_1^{n_1} \bZ)^{\times} \times \cdots \times (\bZ / p_k^{n_k} \bZ)^{\times}} \\
	{\Aut(C_n)} & {\Aut(C_{2^n}) \times \Aut(C_{p_1^{n_1}}) \times \cdots \times \Aut(C_{p_k^{n_k}}) }
	\arrow["\sim", from=1-1, to=1-2]
	\arrow["\sim"', from=1-1, to=2-1]
	\arrow["\sim"', from=1-2, to=2-2]
	\arrow["\sim"', from=2-2, to=2-1]
\end{tikzcd}\] 
and therefore (in the notation of Section \ref{sect:cohomC2modules}) an isomorphism
\[
\frac{(C_n)^0}{(\epsilon + 1)C_n} \xrightarrow{\sim} \frac{(C_{2^r})^0}{(\epsilon + 1)C_{2^r}} \times \frac{(C_{p_1^{n_1}})^0}{(\epsilon + 1)C_{p_1^{n_1}}} \times \cdots \times \frac{(C_{p_k^{n_k}})^0}{(\epsilon + 1)C_{p_k^{n_k}}}.
\]
Then the classification follows from Examples \ref{eg:RealStructuresOddPrimePower} and \ref{eg:RealStructuresEvenPrimePower}, which describe the right-hand side. 
\end{proof}

\subsection{Groups of Odd Order}

\begin{lemma}\label{lem:oddorderdegree2automorphisms}
  Suppose that $H$ has odd order. Then the natural map
  \[
    \Aut(H)[2] \rightarrow \Out(H)[2]
  \]
  is an isomorphism.
\end{lemma}

\begin{proof}
  For injectivity, if $\alpha \in \Aut(H)[2]$ is inner, then $\Ad(g) = \alpha$ for some $g \in H$ and therefore $\alpha$ has odd order. We also have that $\alpha^2 = 1$, hence $\alpha = 1$. For surjectivity, suppose that $\theta^2 = 1$ in $\Out(H)$, and that $\alpha \in \Aut(H)$ is any lift of $\theta$. Then $\alpha^2 = \Ad(g)$ for some $g \in H$, and hence $\alpha^{2n} = 1$ for some odd $n \geq 1$. Then $\alpha^n$ maps to $\theta^n$, and $\theta^n = \theta$ as $n$ is odd, and $(\alpha^n)^2 = 1$, so $\alpha^n \in \Aut(H)[2]$.
\end{proof}

\begin{cor}\label{cor:oddorderRealstructures}
  Suppose that $H$ has odd order. Then
  \[
    \Aut(H)[2] \rightarrow \fR(H), \qquad \alpha \mapsto G(\alpha,1)
  \]
  is a bijection.
\end{cor}

\begin{proof}
  From Lemma \ref{lem:oddorderdegree2automorphisms} we have that the Real structures are partitioned by $\alpha \in \Aut(H)[2]$ rather than $\theta \in \Out(H)[2]$. As $\alpha^2 = 1$, we may always take $h = 1$ and form the semidirect product $G(\alpha, 1)$. By Lemma \ref{cohom_div_by2}, $\HH^2(C_2, Z_{\theta}) = 0$, hence this is the unique Real structure inducing this outer automorphism.
\end{proof}

In particular, we see that all Real structures
\[
1 \rightarrow H \rightarrow G \rightarrow C_2 \rightarrow 1
\]
are split when $H$ has odd order, and the Real structures are exactly $G(\alpha, 1)$ for $\Aut(H)[2]$.

We can upgrade this to the case when $H = H_1 \times H_2$, with $H_2$ of odd order and $\gcd(|H_1|,|H_2|) = 1$. In this case, for $\theta \in \Out(H)[2]$, we write $(\theta_1, \alpha_2)$ for the image of $\theta$ under the isomorphism
\[
\Out(H)[2] \xrightarrow{\sim} \Out(H_1)[2] \times \Out(H_2)[2] \xrightarrow{\sim} \Out(H_1)[2] \times \Aut(H_2)[2]
\]
induced from Lemma \ref{lem:oddorderdegree2automorphisms}.
\begin{cor}\label{cor:Realstructuresoddorderfactor}
  Suppose that $H = H_1 \times H_2$, with $H_2$ of odd order and $\gcd(|H_1|,|H_2|) = 1$. Then for any $\theta \in \Out(H)[2]$, $\fR_{\theta}(H) \neq \emptyset$ if and only if $\fR_{\theta_1}(H_1) \neq \emptyset$, in which case the natural map
  \begin{equation}\label{eqn:oddorderfactorisom}
  \fR_{\theta_1}(H_1) \rightarrow \fR_{\theta}(H), \qquad [G(\alpha_1, h)] \mapsto [G(\alpha_1 \times \alpha_2, (h,1))]
  \end{equation}
  is an isomorphism. In particular, this induces an isomorphism
  \[
    \sR(H_1) \times \Aut(H_2)[2] \xrightarrow{\sim} \sR(H).
  \]
\end{cor}

\begin{remark}
  The surjectivity of the isomorphism (\ref{eqn:oddorderfactorisom}) shows that for any Real structure $G$ on $H$, there is an element $w \in G \setminus H$ with $w^2 \in H_1$. In particular, this allows one to define an explicit inverse to (\ref{eqn:oddorderfactorisom}), by sending $G$ to the subgroup of $G$ generated by $H_1$ and any such $w \in G \setminus H$, which as $w^2 \in H_1$ will be a Real structure on $H_1$.
\end{remark}

\begin{proof}
  First, fix a lift $\alpha_1 \in \Aut(H_1)$ of $\theta_1 \in \Out(H_1)[2]$, and note that $\fR_{\theta_1}(H_1) \neq \emptyset$ if and only if there is some $h \in H_1$ with $\alpha_1(h) = h$ and $\alpha_1^2 = \Ad(h)$ by Lemma \ref{lem:elementary1}. This is equivalent to the existence of a pair $(h,h_2) \in H_1 \times H_2$ with $(\alpha_1(h), \alpha_2(h_2))= (h, h_2)$ and $(\alpha_1 \times \alpha_2)^2 = \Ad(h, h_2)$, i.e.\ that $\fR_{\theta}(H) \neq \emptyset$, as we can always take $h_2 = 1$ because $\alpha_2^2 = 1$.

  Is it straightforward to show that (\ref{eqn:oddorderfactorisom}) is well-defined an injective using Lemma \ref{lem:elementary1} and Lemma \ref{lem:elementary3}. For the surjectivity, any Real structure in $\fR_{\theta}(H)$ has the form $G(\alpha_1 \times \alpha_2, (h,h_2))$ for some $(h, h_2) \in H_1 \times H_2$ and lift $\alpha_1 \in \Aut(H_1)$ of $\theta_1 \in \Out(H_1)[2]$, by Lemma \ref{lem:elementary1} and Lemma \ref{lem:elementary3}. Now $\alpha_2(h_2) = h_2$ and $\alpha_2^2 = \Ad(h_2)$, and so $G(\alpha_2, h_2)$ is a Real structure on $H_2$. By Corollary \ref{cor:oddorderRealstructures}, $G(\alpha_2, h_2) \cong G(\alpha_2, 1)$, and therefore there is some $z \in Z(H_2)$ with $h_2 = z \alpha_2(z)$. In particular, taking $s = (1,z^{-1})$ in Lemma \ref{lem:elementary3}, we see that
\[
G(\alpha_1 \times \alpha_2, (h,h_2)) \cong G(\alpha_1 \times \alpha_2, (h,1))
\]
which shows that (\ref{eqn:oddorderfactorisom}) is surjective.
\end{proof}

\subsection{Centreless Groups} Immediately from Theorem \ref{thm:EMLv1}, we have the following.

\begin{cor}\label{cor:Realstructuresnocenter}
Suppose that $Z(H) = 1$. Then the coupling map
\[
    \fR(H) \rightarrow \Out(H)[2]
\]
is an isomorphism.
\end{cor}

For example, this applies to all non-abelian simple groups. Let us now examine what this means for the symmetric and alternating groups.

\begin{eg}
Suppose that $H=S_n$, for $n \geq 3$, so $Z(S_n) = 1$. When $n \neq 6$, $\Out(S_n) = 1$, and therefore the only Real structure is the trivial Real structure $S_n \leq S_n \times C_2$. When $n = 6$, however, there is an exotic outer automorphism making $|\Out(S_6)| = 2$, and so a second Real structure. This is explicitly given by the inclusion $S_6 \hookrightarrow \Aut(S_6)$.
\end{eg}

\begin{eg}
Suppose that $H=A_n$, for $n \geq 4$, so $Z(A_n) = 1$. When $n \neq 6$, $\Out(A_n)$ has order two, and therefore there are two Real structures on $A_6$ given by $A_n \leq A_n \times C_2$ and $A_n \leq S_n$. When $n = 6$, however, $\Out(A_6) \cong C_2 \times C_2$, and there are four Real structures. In addition to the previous two, there are two other groups which contain $A_6$ as an index two subgroup, namely $\text{PGL}_2(\bF_9)$ and the Mathieu group $M_{10}$.
\end{eg}

\begin{remark}
  When $Z(H) = 1$, from Corollary \ref{cor:Realstructuresnocenter} we know that given $\theta \in \Out(H)[2]$ there is a unique Real structure with coupling $\theta$. If $\alpha \in \Aut(H)$ is a lift of $\theta$, then by Lemma \ref{lem:elementary1} this can be represented by $G(\alpha, h)$ for some $h \in H$. From Lemma \ref{lem:elementary3}, if $G(\alpha, h_1) \cong G(\alpha, h_2)$ for some $h_1, h_2 \in H$, then there is some $s \in H$ with $\Ad(s) = 1$ and $s\alpha(s)h_1 = h_2$, and thus $s = 1$ and $h_1 = h_2$, because $Z(H) = 1$. This shows that for any $\alpha \in \Aut(H)$ which lifts $\theta$, there is a unique $h \in H$ such that $G(\alpha,h)$ forms a group. From the construction of $[\theta] \in \HH^3(C_2, Z_{\theta})$ and the proof of Lemma \ref{lem:deg3cohom}, we can see that this $h$ is the unique $h \in H$ such that $\Ad(h) = \alpha^2$, using again that $Z(H) = 1$.

  In particular, this gives an explicit inverse to the isomorphism of Corollary \ref{cor:Realstructuresnocenter}.
\end{remark}

\section{Weak-Equivalence}\label{sect:weakequiv}

Let $H$ be a group. In this section, we describe how to use the classification of Real structures on $H$ up to isomorphism to classify Real structures on $H$ up to weak equivalence. Here a weak equivalence from $(G_1, H)$ to $(G_2,H)$ is an group isomorphism $\psi \colon G_1 \rightarrow G_2$ for which $\psi(H) = H$. An isomorphism has the additional requirement that $\psi|_H \in \Aut(H)$ is an inner automorphism. We write $G_1 \cong G_2$ (resp.\ $G_1 \simeq G_2$) to denote that there is an isomorphism (resp.\ a weak equivalence) from $G_1$ to $G_2$.

Recall that we write 
\[
\widetilde{\fR}(H) \coloneqq \{ (\alpha,h)\in \Aut(H)\times H \,\mid\, \mbox{equations }\eqref{Conditions_alpha_h}\mbox{ hold}\},
\]
and view $\fR(H)$ as a quotient of this set. The following defines an action of $\Out(H)$ on $\fR(H)$, and shows that the orbits exactly the weak equivalence classes of Real structures on $H$.

\begin{prop}
  Let $H$ be a group. Then:
  \begin{enumerate}
    \item There is an action of $\Aut(H)$ on $\widetilde{\fR}(H)$ defined by setting
    \[
    \beta * (\alpha, h) \coloneqq (\beta  \alpha  \beta^{-1}, \beta(h))
    \]
    for $\beta \in \Aut(H)$ and $(\alpha, h) \in \widetilde{\fR}(H)$. 

  \item If $G(\alpha_1, h_1) \cong G(\alpha_2, h_2)$, then
  \[
  G(\beta  \alpha_1  \beta^{-1} , \beta(h_1)) \cong G(\beta  \alpha_2  \beta^{-1}, \beta(h_2)),
  \]
  and therefore this action induces an action of $\Aut(H)$ on $\fR(H)$.
  
  \item Two equivalence classes $[(\alpha_1,h_1)], [(\alpha_2, h_2)] \in \fR(H)$ are in the same orbit of $\Aut(H)$ if and only if $G(\alpha_1, h_1) \simeq G(\alpha_2, h_2)$.
  
  \item The induced action of $\Aut(H)$ on $\fR(H)$ is trivial on $\Inn(H)$, and thus induces an action of $\Out(H)$ on $\fR(H)$, for which the orbits correspond to weak equivalence classes of Real structures on $H$.
  \end{enumerate}
\end{prop}

\begin{proof}
  For point $(1)$, note that this is a well-defined action: if $\alpha^2 = \Ad(h)$ and $\alpha(h) = h$, then
  \[
  (\beta  \alpha  \beta^{-1})^2 = \beta  \alpha^2  \beta^{-1} = \beta  \Ad(h)  \beta^{-1} = \Ad(\beta(h))
  \]
  and $(\beta  \alpha  \beta^{-1})(\beta(h)) = \beta(h)$, so $\beta * (\alpha, h) \in \widetilde{\fR}(H)$.
  
  For point (2), if $G(\alpha_1, h_1) \cong G(\alpha_2, h_2)$ then there is some $s \in H$ with $s \alpha_2(s) h_2 = h_1$ and $\Ad(s) \alpha_2 = \alpha_1$ by Lemma \ref{lem:elementary3}. Then, we see that 
  \[
  \beta(s) (\beta \alpha_2 \beta^{-1})(\beta(s)) \beta(h_2) = \beta(h_1), \: \text{ and } \: \Ad(\beta(s)) (\beta\alpha_2\beta^{-1}) = \beta\alpha_1\beta^{-1}
  \]
  and therefore $G(\beta  \alpha_1  \beta^{-1} , \beta(h_1)) \cong G(\beta  \alpha_2  \beta^{-1}, \beta(h_2))$, again by Lemma \ref{lem:elementary3}.

  For point (3), suppose that $(\alpha, h) \in \widetilde{\fR}(H)$, and $\beta \in \Aut(H)$. Letting $w$ and $w'$ denote the formal symbols defining $G(\alpha,h)$ and $G(\beta \alpha \beta^{-1}, \beta(h))$ respectively, we define a map
  \[
  \psi \colon G(\alpha,h) \rightarrow G(\beta \alpha \beta^{-1}, \beta(h)) \: \text{ by } \: \psi(h) = \beta(h), \psi(hw) = \beta(h) w',
  \]
  which is clearly bijective by definition. Furthermore, we may check that $\psi(w)^2 = (w')^2 = \beta(h) = \psi(w^2)$, and
  \[
  \psi(wx) = \psi(\alpha(x) w) = \beta(\alpha(x))w' = w' \beta(x) = \psi(w) \psi(x)
  \]
  and therefore this is a well-defined group homomorphism by Corollary \ref{cor:elementary}.

  Conversely, suppose that there is an weak equivalence $\phi \colon G(\alpha_1,h_1) \rightarrow G(\alpha_2, h_2)$. We set $\beta \coloneqq \phi|_{H} \in \Aut(H)$, and claim that
  \[
  G(\alpha_2, h_2) \cong \beta * G(\alpha_1,h_1) = G(\beta \alpha_1 \beta^{-1},\beta(h_1)).
  \]
  By Lemma \ref{lem:elementary3}, this is equivalent to showing that there is some $s \in H$ with
  \[
    s \alpha_2(s) h_2 = \beta(h_1) \text{ and } \Ad(s) \alpha_2 = \beta \alpha_1 \beta^{-1}. 
  \]
  For this, set $s \coloneqq \phi(w_1)w_2^{-1} \in H$, so $\phi(w_1) = s w_2$. Then we may check that
  \[
    \beta(h_1) = \phi(w_1^2) = \phi(w_1)^2 = (s w_2)^2 = s \alpha_2(s)h_2
  \] 
  and that for any $x \in H$,
  \[
  \beta(\alpha_1(\beta^{-1}(x))) s w_2 = \phi(\alpha_1(\beta^{-1}(x)) w_1) = \phi(w_1 \beta^{-1}(x)) = \phi(w_1)\phi(\beta^{-1}(x)) = s w_2 x = s \alpha_2(x) w_2
  \]
  and therefore $\Ad(s) \alpha_2 = \beta \alpha_1 \beta^{-1}$.

  Finally, for point (4), note that for $\beta \in \Aut(H)$, the weak equivalence $\psi \colon G(\alpha,h) \rightarrow G(\beta \alpha \beta^{-1}, \beta(h))$ defined above is actually an isomorphism of Real structures if $\beta \in \Inn(H)$, as $\psi|_H = \beta$. Therefore the action of $\Aut(H)$ on $\Out(H)$ is trivial on $\Inn(H)$.
\end{proof}

\begin{eg}
The smallest example where $\Out(H) \backslash \fR(H)$ is a proper quotient of $\fR(H)$, so there are Real structures which are weak equivalent but not isomorphic, is when $H = C_2 \times C_2 = \langle x \rangle \times \langle y \rangle$. 

The outer automorphism group is
\[
\Out(C_2 \times C_2) = \Aut(\{x,y,xy\}) = S_3,
\]
and therefore the order two elements are $1, (12), (13), (23)$. As in Section \ref{sect:realstrabgroups}, because $C_2 \times C_2$ is abelian we have that for any $\alpha \in \Out(C_2 \times C_2)[2]$ the map 
\[
\frac{\{z \in C_2 \times C_2 \mid \alpha(z) = z\}}{\{z\alpha(z) \mid z \in C_2 \times C_2\}} \rightarrow \fR_{\alpha}(H), \qquad z \mapsto G(\alpha, z)
\]
is an isomorphism. In particular, if $\alpha \neq 1$, then $|\fR_{\alpha}(H)| = 1$. The action of $\Out(C_2 \times C_2)$ on $\Out(C_2 \times C_2)[2]$ by conjugation has two orbits, $\{1\}$ and $\{(12), (13), (23)\}$, and therefore all three Real structures with non-trivial coupling are weak equivalent.

When $\alpha = 1$ instead, $|\fR_{1}(H)| = 4$, and $\fR_{1}(H)$ is explicitly given by 
\[
G(1, 1), \: G(1,x), \: G(1,y), \: G(1,xy).
\]
The first is isomorphic as a group to $C_2^3$, and the remaining three to $C_2 \times C_4$. The stabiliser of $\alpha$ is now $S_3$, and the induced action of $S_3$ on $\fR_{1}(H)$ fixes $G(1,1)$ and identifies all remaining elements.

We therefore see that $H = C_2 \times C_2$ provides examples of distinct Real structures with different couplings which are weak equivalent, and distinct Real structures with the same coupling that are weak equivalent. Overall, $|\fR(H)| = 7$, whereas $|\Out(H) \backslash \fR(H)| = 3$. 
\end{eg}

We now consider how this action of $\Out(H)$ interacts with the decomposition \eqref{RealStr:RTheta} of $\fR(H)$ by coupling, which we use to compute $\Out(H) \backslash \fR(H)$ in practice.

\begin{lemma} \label{lemma:weak_eq}
  For $\eta \in \Out(H)$, $\eta(\fR_{\theta}(H)) = \fR_{\eta \theta \eta^{-1}}(H)$. Furthermore, for any representatives
  \[
  \{\theta_t \mid t \in \Out(H) \backslash \Out(H)[2]\} \subset \Out(H)[2]
  \]
  of $\Out(H) \backslash \Out(H)[2]$, where $\Out(H)$ acts on $\Out(H)[2]$ by conjugation, the natural map
  \[
  \bigsqcup_{t \in \Out(H) \backslash \Out(H)[2]} \Out_{\theta_t}(H) \backslash \fR_{\theta_t}(H) \longrightarrow \Out(H) \backslash \fR(H)
  \]
  is an isomorphism, where $\Out_{\theta_t}(H) \coloneqq C_{\Out(H)}(\theta_t)$ denotes the centraliser of $\theta_t$ in $\Out(H)$.
\end{lemma}

\begin{proof}
  For any $\theta, \eta \in \Out(H)$, any element of $\fR_{\theta}(H)$ is represented by some $G(\alpha, h)$ with $[\alpha] = \theta$ in $\Out(H)$, and therefore $\eta * G(\alpha, h)$ has coupling $\eta [\alpha] \eta^{-1} = \eta \theta \eta^{-1}$.

  The second claim follows from the following general fact. Suppose that $X$ is a left $G$-set, and $X = \sqcup_{i \in I} X_i$ is a disjoint union decomposition which is permuted by $G$, meaning that for any $g \in G$ and $i \in I$, $g(X_i) = X_{g(i)}$ for some $g(i) \in I$. Then this defines an action of $G$ on $I$, and for any choice of representatives $\{i_t \mid t \in G \backslash I\} \subset I$ of $G \backslash I$, the natural map
  \[
  \bigsqcup_{t \in G \setminus I} \Stab_{G}(i_t) \backslash X_{i_t} \rightarrow G \backslash X
  \]
  is a bijection.
\end{proof}

Given $\theta \in \Out(H)[2]$, if $\fR_{\theta}(H)$ is non-empty, then we can describe $\Out_{\theta}(H) \backslash \fR_{\theta}(H)$ explicitly as follows. Fix a lift $\alpha \in \Aut(H)$ of $\theta$, and let $h \in H$ be such that $G(\alpha,h)$ forms a group, or equivalently that $\Ad(h) = \alpha^2$ and $\alpha(h) = h$. For example, this may always be done by taking $t \in H$ with $\Ad(t) = \alpha^2$, finding $z \in Z(H)$ with $\alpha(t)t^{-1} = \alpha(z) z^{-1}$ and setting $h \coloneqq tz^{-1}$, as in Lemma \ref{lem:deg3cohom}.

From Lemma \ref{lem:freetransaction}, we see that the map
\[
\Phi_{\alpha, h} \colon Q_{\theta} \coloneqq Z_{\theta}^{C_2} / (\epsilon + 1)Z_{\theta} \rightarrow \fR_{\theta}(H), \qquad \Phi_{\alpha, h}(z) \coloneqq G(\alpha, z h)
\]
is a bijection. The following lemma describes an action of $\Out_{\theta}(H)$ on $Q_{\theta}$ which makes this bijection equivariant with respect to the action described above of $\Out_{\theta}(H)$ on $\fR_{\theta}(H)$.

\begin{lemma}\label{lem:torsoraction}
Let $\alpha\in\Aut(H)$ and $h\in H$ satisfy $\alpha^2=\Ad(h)$ and $\alpha(h)=h$. Let $\beta\in\Aut(H)$ satisfy
$[\alpha,\beta]=1$ in $\Out(H)$ and choose $s\in H$ such that $\beta\alpha\beta^{-1}=\Ad(s)\alpha$. Set
\[
z_\beta\coloneqq
(s\alpha(s))^{-1}\beta(h)h^{-1}.
\]
Then $z_\beta\in Z_\theta^{C_2}$, its class in $Q_{\theta}$ is independent of the choice of $s$, and
\[
[\beta]\ast[z]\coloneqq[\beta(z)z_\beta]
\]
defines a well-defined action of $\Out_\theta(H)$ on $Q_{\theta}$ making $\Phi_{\alpha,h}$ $\Out_\theta(H)$-equivariant.
\end{lemma}

\begin{proof}
  It is possible to check this directly by hand, but we instead proceed as follows, and first compute the induced action of $\Out_\theta(H)$ on $Q_{\theta}$ transported via $\Phi_{\alpha,h}$. For $[z] \in Q_{\theta}$, and $[\beta] \in \Out_{\theta}(H)$, 
  \[
  [\beta] * \Phi_{\alpha, h}([z]) = [\beta] * G(\alpha, zh) = G(\beta \alpha \beta^{-1}, \beta(z) \beta(h)),
  \]
  and therefore $[\beta] * [z] \in Q_{\theta}$ is the unique element of $Q_{\theta}$ such that there is some $s \in H$ with $\beta \alpha \beta^{-1} = \Ad(s)$ and $s \alpha(s) ([\beta] * [z]) h = \beta(z) \beta(h)$, by Lemma \ref{lem:elementary3}. Note that if $s' \in H$ is any element with $\beta \alpha \beta^{-1} = \Ad(s)$, then $s' = sy$ for some $y \in Z(H)$, and therefore
  \[
  (s' \alpha(s'))^{-1}\beta(z) \beta(h)h^{-1} = (y \alpha(y))^{-1} (s \alpha(s))^{-1}\beta(z) \beta(h)h^{-1} = (s \alpha(s))^{-1}\beta(z) \beta(h)h^{-1}
  \]
  in $Q_{\theta}$, and therefore 
  \[
  [\beta] * [z] = \beta(z) (s \alpha(s))^{-1} \beta(h) h^{-1} = \beta(z) z_{\beta}
  \]
  in $Q_{\theta}$ for \emph{any} $s \in H$ with $\beta \alpha \beta^{-1} = \Ad(s)$. As $[\beta] * [z] \in Q_{\theta}$, this tells us that $\beta(z) z_{\beta}$ and thus $z_{\beta}$ are in $Z_{\theta}^{C_2}$, $[z_{\beta}] \in Q_{\theta}$ is independent of $s$, and defining $[\beta]\ast[z]\coloneqq[\beta(z)z_\beta]$ gives exactly the transported action of $\Out_\theta(H)$ on $Q_{\theta}$ under $\Phi_{\alpha,h}$, as required.
\end{proof}

\subsection{Algorithm} In \cite{GITHUB}, we give an algorithm in Magma which, for a finite group $H$, returns all possible Real structures on $H$ up to isomorphism, and up to weak equivalence.

Below we describe this algorithm, which is based in the results of Section \ref{sect:unification} and Section \ref{sect:weakequiv}. This applies to any group, finite or not, for which each subroutine of the algorithm can be implemented.

\begin{algorithm}
\caption{Compute $\fR_\theta(H)$ and $\Out_{\theta}(H) \backslash \fR_{\theta}(H)$ for an order two element $\theta \in \mathrm{Out}(H)$}
\begin{algorithmic}[1]

\Require A group $H$, an order two element $\theta \in \mathrm{Out}(H)$, an element $\alpha \in \mathrm{Aut}(H)$ lifting $\theta$
\Ensure A list of elements of $H$, $\mathfrak{h}$, such that $\{G(\alpha, h) : h \in \mathfrak{h}\}$ is a set of representatives for $\fR_\theta(H)$, and a subset $\mathfrak{h}_0$ such that $\{G(\alpha, h) : h \in \mathfrak{h}_0\}$ is a set of representatives for $\Out_{\theta}(H) \backslash \fR_{\theta}(H)$

    \State Find $t \in H$ such that $\alpha^2 = \mathrm{Ad}(t)$
    \State Set $\lambda := \alpha(t) t^{-1} \in H$
    \If{there exists $z \in Z(H)$ such that $\lambda = \alpha(z) z^{-1}$}
        \State Fix such a $z$ and set $h := tz^{-1}$
        \State Compute the subgroups of $Z(H)$:
        \[
            Z^\alpha := \{ z \in Z(H) : \alpha(z) = z \}
        \]
        \[
            B^\alpha := \{ \alpha(z) \cdot z : z \in Z(H) \}
        \]
        \State Set $Q \coloneqq Z^\alpha / B^\alpha, \pi \colon Z^{\alpha} \twoheadrightarrow Q$ the quotient map, $\sigma \colon Q \rightarrow Z^{\alpha}$ a section
        \State Let $\mathfrak{B}$ be a set of coset representatives in $\Aut(H)$ of $\eta \in \Out_{\theta}(H)$
        \State For each $\beta \in \mathfrak{B}$, find $s \in H$ with $\beta \alpha \beta^{-1} = \Ad(s) \alpha$, and set
        \[
            z_{\beta} \coloneqq (s \alpha(s))^{-1}\beta(h) h^{-1}
        \]
        \State Set $\mathrm{orbitreps} \leftarrow []$, $\mathrm{tosort} \leftarrow Q$
        \While{$\mathrm{tosort} \neq \emptyset$}
          \State Pick $z \in \mathrm{tosort}$
          \State $\mathrm{orbit} \coloneqq \{\beta(z)z_{\beta} \mid \beta \in \mathfrak{B}\} \subset Q$
          \State Add $z$ to $\mathrm{orbitreps}$
          \State Remove $\mathrm{orbit}$ from $\mathrm{tosort}$
        \EndWhile

        \State $\mathfrak{h} \leftarrow [\sigma(z)h : z \in Q]$
        \State $\mathfrak{h}_0 \leftarrow [\sigma(z)h : z \in \mathrm{orbitreps}]$
    \Else
        \State $\mathfrak{h} \leftarrow [\,]$
        \State $\mathfrak{h}_0 \leftarrow [\,]$
    \EndIf

\Return $\mathfrak{h}, \mathfrak{h}_0$

\end{algorithmic}
\end{algorithm}

\begin{algorithm}
\caption{Compute $\fR(H)$ and $\Out(H) \backslash \fR(H)$}
\begin{algorithmic}[1]

\Require A group $H$
\Ensure A list of automorphisms $\mathfrak{a}$ of $H$, which represent $\Out(H)[2]$ in $\Aut(H)$, and a subset $\mathfrak{a}_0 \subset \mathfrak{a}$ which whose reductions in $\Out(H)[2]$ represent $\Out(H) \backslash \Out(H)[2]$. For each $\alpha \in \mathfrak{a}$, a list of elements $\mathfrak{h}_{\alpha}$ of $H$, such that $\{G(\alpha, h) : h \in \mathfrak{h}_{\alpha}\}$ is a set of representatives for $\fR_{[\alpha]}(H)$. For each $\alpha \in \mathfrak{a}_0$, a list of elements $\mathfrak{h}_{\alpha, 0}$ of $H$, such that $\{G(\alpha, h) : h \in \mathfrak{h}_{\alpha, 0}\}$ is a set of representatives for $\Out_{[\alpha]}(H) \backslash \fR_{[\alpha]}(H)$.

In particular,
\[
\{G(\alpha, h) : \alpha \in \mathfrak{a}, h \in \mathfrak{h}_{\alpha}\}
\]
is a set of representatives of $\fR(H)$, and 
\[
\{G(\alpha, h) : \alpha \in \mathfrak{a}_0, h \in \mathfrak{h}_{\alpha,0}\}
\]
is a set of representatives for $\Out(H) \backslash \fR(H)$.

\State Compute $\Out(H)[2]$
\State Pick a section $s \colon \Out(H) \rightarrow \Aut(H)$
\State Pick $\{\theta_t \colon t \in \Out(H) \backslash \Out(H)[2]\}$, representatives for $\Out(H) \backslash \Out(H)[2]$
\State $\mathfrak{a} \leftarrow [s(\theta) \mid \theta \in \Out(H)[2]]$
\State $\mathfrak{a}_0 \leftarrow [s(\theta_t) \mid t \in \Out(H) \backslash \Out(H)[2]]$

\State For each $\alpha \in \mathfrak{a}$, $\mathfrak{h}_{\alpha}$ is the $\mathfrak{h}$ provided by Algorithm 1
\State For each $\alpha \in \mathfrak{a}_0$, $\mathfrak{h}_{\alpha, 0}$ is the $\mathfrak{h}_0$ provided by Algorithm 1

\Return $\mathfrak{a}, \mathfrak{a}_0, [\mathfrak{h}_{\alpha} \mid \alpha \in \mathfrak{a}], [\mathfrak{h}_{\alpha, 0} \mid \alpha \in \mathfrak{a}_0]$

\end{algorithmic}
\end{algorithm}
